\section{Experimental Setup}\label{sec:setup}
TODO: Describe the technical implementation (e.g., framework used, LLM endpoints).

\subsection{Quantitative Scoring System}\label{subsec:scoring}

\subsubsection{Biologist Agent Scoring}
The objective is to rank a batch of peptides and identify which candidate is closest to the provided context (an existing peptide or requested properties represented as a reference embedding).

The implementation uses ESM-2 embeddings (default: \texttt{facebook/esm2\_t12\_35M\_UR50D}) with attention-mask-aware mean pooling to produce one vector per peptide. We first tested cosine similarity between candidate and reference embeddings. In practice, cosine values were often too similar for short peptides in this protein-wide latent space, which reduced ranking contrast. We therefore switched to Euclidean distance, then normalized batch distances with a Z-score and mapped them to $[0,1]$ with a sigmoid.

The resulting score is used for relative ranking inside each batch and for Top-$K$ selection by the orchestrator. It should be interpreted as a proxy signal to guide the search process, not as an absolute biological activity measurement.

\subsubsection{Chemist Agent Scoring}

Crucially, to prevent pipeline starvation (where strict filters might discard the entire population early in the optimization process), the agent employs an adaptive thresholding mechanism. The workflow operates in two phases:
\begin{enumerate}
    \item \textbf{Hard Constraint Check:} Candidates are first filtered against strict limits (e.g., range constraints on length or charge).
    \item \textbf{Population Rescue:} If the survival rate falls below a predefined threshold (e.g., $50\%$), the agent relaxes the hard constraints. It then fills the remaining quota by selecting the top-ranking rejected candidates—those that, while slightly violating limits, are mathematically closest to the ideal TPP.
\end{enumerate}
This strategy ensures a continuous flow of candidates to the Biology Agent while maximizing selective pressure towards developable molecules.

\begin{itemize}
    \item \textbf{Solubility \& Stability:} We calculate the Isoelectric Point (pI) and Net Charge at specific physiological pH levels (e.g., pH 5.5). These metrics are critical for predicting aggregation propensity and solubility in biological fluids.
    \item \textbf{Bioavailability \& Permeability:} The generated profile includes the partition coefficient ($\log P$), Topological Polar Surface Area (TPSA), and hydrophobicity indices. These values act as proxies for the peptide's ability to cross biological membranes and interact with hydrophobic pockets.
    \item \textbf{Structural Indices:} Metrics such as the Boman Index are computed to estimate potential protein-binding potential or antimicrobial activity.
\end{itemize}

\subsubsection{Orchestrator-Level Scoring and Ranking}
TODO: Describe how the orchestrator combines outputs from expert agents (Biologist + Chemist), applies weighting or filtering rules, and computes the final ranking used for Top-$K$ selection.

\subsection{Evaluation Metrics}\label{subsec:metrics}
TODO: Define the metrics you tracked: Validity (\%), Diversity (distances/entropy/n-grams), and Constraint pass rates[cite: 25, 27, 28].

\subsection{Baselines}\label{subsec:baselines}
The baseline is designed as a simpler non-agentic reference to test whether the proposed multi-expert iterative pipeline provides measurable gains. Concretely, we use a standalone Conditional Variational Autoencoder (CVAE) that directly generates peptide candidates from conditioning features, without chemistry/bio critique loops and without orchestrated revision rounds.

This baseline CVAE is trained on the dataset built with \texttt{database/get\_data.py}. The extraction script queries DBAASP asynchronously and retains sequence-level and physicochemical fields used for conditioning (length, pH, molecular weight, logP, net charge, isoelectric point, hydrophobicity, and cationicity). In practice, many API entries contain missing values for one or more properties. The baseline data pipeline therefore substitutes absent fields with fixed defaults before standardization. This dataset is thus not optimal from a data-quality perspective, but for a proof-of-concept objective we considered it sufficient to establish a consistent comparison point.

From an implementation perspective, the baseline model is a GRU-based CVAE. The encoder embeds amino-acid tokens and maps sequence-condition pairs to latent Gaussian parameters; the decoder is autoregressive and conditioned on both latent code and property vector. Training uses cross-entropy reconstruction loss with KL regularization, Adam optimization, and a progressive KL weight (annealed $\beta$).

At inference time, candidate peptides are sampled from the latent prior under a target condition vector (scaled with the training \texttt{StandardScaler}), with temperature and top-$k$ sampling to control exploration. This baseline serves only as the comparison anchor: unlike the full pipeline, it performs one-shot conditional generation and does not benefit from multi-agent feedback, constraint-aware iteration, or round-level refinement.

\begin{figure*}[t]
\centering
\begin{tikzpicture}[
	>=latex,
	modelBlock/.style={
		rectangle,
		thick,
		minimum width=3cm,
		minimum height=1.5cm,
		align=center,
		font=\sffamily\bfseries,
		rounded corners=2pt
	},
	dataBlock/.style={
		rectangle,
		thick, dashed,
		minimum width=2.5cm,
		minimum height=1cm,
		align=center,
		font=\sffamily\footnotesize,
		rounded corners=2pt,
		fill=gray!5, draw=gray!60
	},
	condBlock/.style={
		rectangle,
		thick,
		minimum width=6.6cm,
		minimum height=1.2cm,
		align=center,
		font=\sffamily\footnotesize,
		rounded corners=2pt,
		fill=orange!10, draw=orange!80!black
	},
	encStyle/.style={modelBlock, fill=blue!10, draw=blue!80!black},
	decStyle/.style={modelBlock, fill=green!10, draw=green!80!black},
	latentStyle/.style={modelBlock, fill=purple!10, draw=purple!80!black, minimum width=2.7cm},
	outStyle/.style={dataBlock, solid, fill=yellow!15, draw=yellow!60!black},
	flowArrow/.style={-{Stealth[length=3mm]}, thick, draw=black!90},
	condArrow/.style={-{Stealth[length=3mm]}, thick, draw=orange!90!black},
	inferArrow/.style={-{Stealth[length=3mm]}, thick, dotted, draw=purple!80!black}
]

\node[dataBlock] (inputSeq) at (0, 0) {Input Sequence $x$\\ \scriptsize(Training)};
\node[encStyle] (encoder) at (4, 0) {Encoder\\ \scriptsize(Embedding + GRU)};
\node[latentStyle] (latent) at (8, 0) {Latent \\ \scriptsize$z=\mu+\sigma\odot\epsilon$};
\node[decStyle] (decoder) at (12, 0) {Decoder\\ \scriptsize(Autoregressive GRU)};
\node[outStyle] (outputSeq) at (16, 0) {Generated\\ peptide $\hat{x}$};

\node[condBlock] (condition) at (8, -2.5) {\textbf{Condition vector} $c$\\ \scriptsize [length, pH, MW, logP, net charge, isoP, hydrophobicity, cationicity]};

\node[dataBlock, draw=purple!60, fill=purple!5] (noise) at (8, 2.5) {Prior sample\\ $\mathcal{N}(0, I)$\\ \scriptsize(Inference)};

\draw[flowArrow] (inputSeq) -- (encoder);
\draw[flowArrow] (encoder) -- (latent);
\draw[flowArrow] (latent) -- (decoder);
\draw[flowArrow] (decoder) -- (outputSeq);

\draw[condArrow] (condition.west) -| node[pos=0.75, left, font=\sffamily\scriptsize, text=orange!80!black] {concat $c$} (encoder.south);
\draw[condArrow] (condition.east) -| node[pos=0.75, right, font=\sffamily\scriptsize, text=orange!80!black] {concat $c$} (decoder.south);

\draw[inferArrow] (noise) -- (latent);

\end{tikzpicture}
\caption{Baseline CVAE architecture used for comparison. During training, the model learns $q_{\phi}(z\mid x,c)$ and reconstructs sequences with an autoregressive decoder conditioned on $z$ and $c$. During inference, peptides are generated by sampling $z\sim\mathcal{N}(0,I)$ and decoding once, without iterative multi-agent feedback.}
\label{fig:baseline_cvae}
\end{figure*}
